
La investigación en inteligencia artificial (IA) comenzó a ganar popularidad a partir de mediados del siglo XX. El término``inteligencia artificial'' fue propuesto en 1956 durante la conferencia de Dartmouth, que se considera un hito fundacional en el campo de la IA \cite{McCarthy_Minsky_Rochester_Shannon_2006}. Sin embargo, la investigación y el interés en esta área han crecido significativamente en las últimas décadas, especialmente en los años 1980 y 1990 con el avance de las tecnologías de computación y la disponibilidad de mayores capacidades de procesamiento \cite{aiindex2023, LeCun2015}.

La optimización es un elemento crítico en diversas industrias, desde la logística hasta la manufactura, donde se busca maximizar la eficiencia y minimizar costos. Varios métodos tradicionales o heurísticos fueron desarrollados para lograr estos objetivos de optimización \cite{CHEN201771, bard1990branch, vince2002framework}. Más tarde, Fred Glover, acuñó el término ``metaheurística'' en 1986 \cite{glover1986future} para describir métodos heurísticos con características no específicas del problema \cite{abdel2018metaheuristic}. Dentro de la IA, se han destacado subcampos como el aprendizaje automático (Machine Learning), aprendizaje profundo (Deep Learning), y la optimización evolutiva (Evolutionary Algorithms). El aprendizaje automático y profundo abarcan técnicas como las redes neuronales artificiales, máquinas de soporte vectorial, redes convolucionales, entre otras. Estas se utilizan para identificar patrones y realizar predicciones a partir de grandes conjuntos de datos. Recientemente, surgió el término Inteligencia Artificial Explicable (XAI) en respuesta a la creciente complejidad y opacidad de los modelos modernos de IA, especialmente en aplicaciones críticas donde es escencial comprender el proceso de toma de decisiones para la confianza y la responsabilidad. XAI busca hacer que los sistemas de IA sean más transparentes e interpretables.

Por otra parte, la optimización evolutiva utiliza algoritmos basados en principios de evolución natural, como los algoritmos genéticos, para encontrar soluciones óptimas en espacios de búsqueda complejos \cite{izzo2019survey}. En el  contexto de la optimización metaheurística, se han desarrollado diversos algoritmos que proponen técnicas prometedoras para solucionar los problemas de optimización. Algunos de estos algoritmos son bioinspirados como  Ant Colony Optimization (ACO) \cite{484436}, Particle Swarm Optimization (PSO) \cite{488968}, Bat Algorithm (BA) \cite{bat-algorithm}, Dragonfly Algorithm (DA) \cite{Meraihi2020}, Slap Swarm Algorithm (SSA) \cite{MIRJALILI2017163}, entre otros que han surgido en las últimas decadas.

Sin embargo, a pesar de los avances en algoritmos de optimización y sus variantes de mejoras e hibridación, persiste la problemática del estancamiento, donde estos algoritmos se ven atrapados en soluciones subóptimas debido a la falta de exploración en el espacio de búsqueda. Esta problemática afecta a los sectores productivos que utilizan algoritmos de optimización para obtener beneficios diversos. Por ejemplo, las empresas de transporte que deben decidir las rutas más rentables o las organizaciones que enfrentan problemas complejos de optimización. En particular, el Problema del Recorrido Rentable (PTP), que combina beneficios y costos de viaje en una función objetivo, pone de manifiesto la necesidad de estrategias más eficaces para superar el estancamiento y mejorar la toma de decisiones.

Observando en la literatura un fuerte avance en la interpretabilidad de modelos por parte de XAI, la propuesta de este trabajo consiste en un enfoque híbrido que integra modelos bayesianos y técnicas de explicabilidad XAI, como SHAP, para abordar el estancamiento en algoritmos de optimización. Esta metodología no solo busca mejorar la calidad de las soluciones encontradas, sino también proporcionar un entendimiento claro de los factores que contribuyen a la generación de estancamientos.

Las contribuciones de este estudio pueden ser significativas, ya que aportarían un nuevo enfoque para la optimización de algoritmos, abriendo nuevas áreas de investigación en metaheurísticas. Además, la aplicación del método propuesto a un problema real como el PTP ofrecerá no solo soluciones más eficientes, sino que también permitirá a las empresas optimizar sus operaciones y mejorar la satisfacción del cliente. Al combinar la teoría con la práctica, este trabajo busca abordar problemas académicos y ofrecer soluciones aplicables y útiles en el mundo real.


La organización de este trabajo de investigación es la siguiente: el Capítulo 2 presenta el marco conceptual \ref{sc:MC} del trabajo junto al estado del arte \ref{sc:EA}. El Capítulo 3 se expone la descripción del problema \ref{sc:FP}, la propuesta de solución \ref{sc:SP}, preguntas de investigación \ref{sc:PI}, hipótesis \ref{sc:HIPO}, los objetivos generales \ref{sc:OG} y específicos \ref{sc:OE}. En el Capítulo 4 \ref{sc:METOD} se detalla la metodología empleada, incluyendo el diseño de las experimentaciones, el problema de optimización a resolver, el análisis de resultados y cómo serán reportados. El Capítulo 5 describe el plan de trabajo \ref{sc:PT}. Finalmente, en el Capítulo 6 se especifican los recursos humanos \ref{sc:RH} y materiales \ref{sc:RM} requeridos durante la investigación.
